\section{Rule and State Checks Support Correct Complete-Game Execution}

The automated suite contains 47 passing tests covering movement rules, opening rolls, outcome deduplication, heuristic symmetry, rollout depth, turn alternation, and scoring perspective. Rule fixtures check both players through mirrored positions. An earlier 100-game random-versus-random run completed without state or termination failures. Together, these results support using the engine for pilot experiments. They do not establish exhaustive correctness or agreement with an independent expert engine.

Table~\ref{tab:baseline-results} summarizes the retained complete-game experiments. The main agent alternates equally between the two player identities within each batch; the opening roll independently determines who starts. Report settings, rather than filenames, identify the configurations.


\begin{table}[htbp]
\centering\small
\caption{Complete-game performance of the evaluated agent against random and heuristic-only opponents. Wins and win rates refer to the evaluated agent. A rollout is one simulated continuation per distinct candidate; depth counts subsequent individual player turns, excluding the candidate, and simulations stop early at a terminal state. Guided moves maximize the acting player's immediate heuristic score; heuristic-only play evaluates candidates without simulating future turns. Each 40-game batch assigns the evaluated agent to each player identity 20 times. Rows pool different seed batches and vary in policy, rollout count, or depth, so the table is an exploratory summary rather than an isolated comparison of any single factor.}
\label{tab:baseline-results}
\resizebox{\linewidth}{!}{%
\begin{tabular}{@{}llrrrl@{}}
\toprule
Opponent & Evaluated policy & Rollouts per candidate & Depth in turns & Wins / games & Win rate \\
\midrule
Random & Random rollouts & 1 & 1 & 159 / 160 & 99.4\% \\
Random & Heuristic-only & -- & 0 & 40 / 40 & 100.0\% \\
Heuristic-only & Random rollouts & 20 & 20 & 4 / 40 & 10.0\% \\
Heuristic-only & Guided rollouts & 20 & 4 & 53 / 80 & 66.3\% \\
Heuristic-only & Guided rollouts & 5 & 4 & 23 / 80 & 28.8\% \\
\bottomrule
\end{tabular}}
\end{table}

The first row pools base seeds 20261017, 5000, 6000, and 7000. The heuristic-only and random-depth-20 rows each use seed 20261017. The guided 20-rollout row pools seeds 20261017 and 5000; the guided five-rollout row pools seeds 10000 and 11000.

\section{Success Against Random Play Does Not Establish the Value of Lookahead}

The shallow random-rollout agent won 159 of 160 games against random play. However, heuristic-only play also won all 40 games in its random-opponent batch. These outcomes show an advantage over random choices but leave little room to distinguish the value of search. Against heuristic-only play, the retained random-rollout configuration with 20 rollouts and depth 20 won only four of 40 games. Thus, success against a weak opponent was insufficient validation of the search method.

A plausible explanation is that random continuations evaluate positions under behavior unlike the actual opponent's. Actions that benefit from random mistakes need not perform well against deliberate heuristic choices. This interpretation concerns the rollout model, not a demonstrated failure of Monte Carlo search in general. The experiments do not isolate the mechanism, and the retained configurations differ in depth and opponent.

\section{Guided Rollouts Show a Promising Advantage Over Immediate Evaluation}

At depth four with 20 rollouts per candidate, guided search won 29/40 games in one retained batch and 24/40 in another. The pooled result is 53/80 wins, or 66.3\%, with an approximate 95\% Wilson interval of 55.4\%-75.7\%. This supports a provisional advantage over the shared immediate heuristic in this matchup, while the variation between batches shows why one batch is insufficient.

Guided rollouts choose sensible immediate moves for each simulated player while retaining stochastic dice and random tie-breaking. They therefore more closely match the heuristic-only opponent than random continuations do. This offers a plausible explanation for the improvement, although the random and guided rows in Table~\ref{tab:baseline-results} also differ in depth. A matched policy comparison is needed before attributing the difference solely to rollout guidance.

This finding motivates validating agent behavior before studying rule variations. It is consistent with the methodological emphasis on agent-based playtesting in de Mesentier Silva et al. \cite{silva2017}, but their Ticket to Ride setting differs from this project. The present win rates are not directly comparable to that study or to AlphaZero's chess-variant results \cite{tomasev2020}. No modified Backgammon rules have been tested yet.

\section{Five-Rollout Results Suggest Sensitivity to Sampling Budget}

With guided rollouts and the same depth of four, reducing the budget to five rollouts produced 12/40 and 11/40 wins, totaling 23/80 or 28.8\%. One possible explanation is noisy action ranking: a small sample can make an inferior action appear favorable, especially when the agent selects the highest estimate among many candidates.

The observed difference from the 20-rollout results is large, but the pooled groups use different seed batches. These experiments therefore suggest budget sensitivity rather than establish a causal effect or an optimal rollout count. The 20-rollout configuration remains the provisional choice; a matched seed comparison and held-out evaluation should precede stronger claims.

\section{Typical Decision Times Are Practical but Occasional Decisions Are Slow}

The retained 20-rollout, depth-four guided batches recorded 3,215 main-agent decisions, including forced passes. Their pooled mean decision time was approximately 0.323 seconds, calculated over decisions rather than averaging the two batch means equally. The individual batch means were 0.347 and 0.299 seconds. Their within-batch 95th-percentile times were 1.273 and 1.148 seconds, while maximum times reached 16.002 and 16.367 seconds.

Typical decisions are fast enough for continued development, but the long tail matters when planning large tournaments. Cost depends on legal candidate counts and the simulated positions encountered. Guided rollouts evaluate candidate positions at each simulated turn, so equal rollout counts across policies do not represent equal computational budgets. These local measurements exclude the evaluator's legality and state-checking overhead; total experiment time is greater.

\section{Preserved Evidence Supports a Provisional Baseline, Not Final Research Conclusions}

The current operational baseline is guided rollouts, depth four, and 20 rollouts per distinct candidate. These settings are neither an established optimum nor evidence of expert playing strength. The heuristic weights remain hand-selected, and the reported intervals are exploratory because the experiments informed configuration selection.

An additional guided 20-rollout seed-10000 batch was reported as 26/40 wins in terminal output, but its JSON was overwritten by a five-rollout run. It is excluded from Table~\ref{tab:baseline-results} and the retained-data aggregate. Two filenames containing "r20" currently record five rollouts; their settings determine how they are classified. Unique experiment identifiers and an experiment ledger are needed before formal data collection.

The midterm results establish an operational simulator and a promising search baseline. They do not yet establish distinct strategic styles or the effects of rule changes. Those questions require behavioral validation of the four planned agents followed by controlled standard-rule and modified-rule tournaments.
